Table~\ref{tab:fp-wm} and Fig.~\ref{fig:fp-wm} combine each modality's MIT-licensed watermark with its fingerprints on the same attacked copies of watermarked assets. The two satisfy their respective success criteria on different queries. For images, TrustMark yields its key for \val{fp:wm/image/DINOv2-S/wm/pct}\% of attacked copies and DINOv2-S identifies \val{fp:wm/image/DINOv2-S/fp/pct}\%; the union of these successes covers \val{fp:wm/image/DINOv2-S/either/pct}\%, and \val{fp:wm/image/ISC21-1st/either/pct}\% with the ISC21 descriptor. The watermark alone survives additive noise, emoji and text overlays, perspective and the re-capture chain better than DINOv2-S; the fingerprint alone survives the meme layout, padding, a 5-degree rotation and skew, which move the watermark out of TrustMark's decoding grid. Neither survives large crops, rotations of 15 degrees or more, screenshots or overlays onto another photograph (Fig.~\ref{fig:fp-wm}a). For audio the fingerprint covers more queries: AudioSeal returned the exact 16-bit key for \val{fp:wm/audio/Chromaprint/wm/pct}\% of queries (\val{fp:wm/audio/keyok/unique/pct}\% if recoveries of the \val{fp:wm/audio/n_colliding} assets that share a key are not counted as identification) and Chromaprint identified \val{fp:wm/audio/Chromaprint/fp/pct}\%, and the watermark added a useful share only under telephone filtering and loud babble; nothing survived pitch, tempo or speed changes. For video expected-key watermark verification covers more queries: Video Seal passed that test for \val{fp:wm/video/ISCC-Video-64/wm/pct}\% of queries, including the crops, flips, rotations, letterboxing and insertions that break every video fingerprint, and the fingerprints added almost nothing (a few screen recordings); picture-in-picture defeated both. These union rates describe complementary coverage, not the measured accuracy of a deployed cascade; the video result does not establish unique registry-ID recovery.

\paragraph{Watermarking moves fingerprints.} Embedding the watermark is itself a transformation. For \val{fp:wm/image/ISCC-Image-64/drift/pct}\% of the watermarked images the ISCC Image-Code of the marked image no longer matched the unmarked original at the operating threshold, and for \val{fp:wm/image/SSCD-mixup/drift/pct}\% with SSCD, against \val{fp:wm/image/PDQ/drift/pct}\% for PDQ and \val{fp:wm/image/DINOv2-S/drift/pct}\% for DINOv2-S. We attribute the ISCC case to the same step that makes it survive padding: on product shots with white backgrounds, a faint residual in the border changes what the uniform-border trim removes, and the security experiments show the same sensitivity to imperceptible perturbations. A registry that stores the fingerprint of the unmarked original therefore loses matches: ISCC's detection of attacked marked copies fell from \val{fp:wm/image/ISCC-Image-64/fp/pct}\% to \val{fp:wm/image/ISCC-Image-64/unmarkedreg/fp/pct}\% when the registry held the unmarked fingerprint. The fingerprint must be computed on the published, watermarked asset.

\paragraph{False keys.} On never-registered, unmarked images TrustMark's error-correcting decoder returned a valid codeword for \val{fp:wm/image/falsekey/pct}\% of \val{fp:wm/image/falsekey/n} attacked queries (a detection; these decodes were not compared with the registered keys), in line with its blind false-positive rates in the watermarking track (Section~\ref{sec:wm-res-fpr}); AudioSeal's detector fired on \val{fp:wm/audio/falsekey/k} of \val{fp:wm/audio/falsekey/n} attacked queries from \val{fp:wm/audio/falsekey/nsrc} never-registered tracks, and that detection's message equalled none of the \val{fp:wm/audio/nkeys} distinct registered keys; and Video Seal, whose key needs agreement on far more bits, matched a registered key for \val{fp:wm/video/falsekey/pct}\%. A cascade must check a decoded key against the registry before trusting it, exactly as it thresholds a fingerprint match.

\begin{table}[!t]
\centering
\footnotesize
\caption{Watermark and fingerprint on the same attacked copies of watermarked assets, pooled over transformations. WM: exact-key recovery for TrustMark and AudioSeal, but expected-key verification for Video Seal; the latter does not establish unique registry identification. FP: correct fingerprint identification at a threshold calibrated to pair-level FPR $10^{-7}$ for images or query-level FPR 1\% for audio and video; achieved held-out rates are reported separately. Both and Either are the intersection and union of these success criteria, not measured cascade accuracy. Drift: share of watermarked assets whose own fingerprint falls below the threshold against the unmarked original.}
\label{tab:fp-wm}
\setlength{\tabcolsep}{3pt}
\begin{tabular}{llccccc}
\toprule
Watermark & Fingerprint & WM & FP & Both & Either & Drift \\
\midrule
TrustMark-Q & DINOv2-S & 0.62 & 0.61 & 0.55 & 0.69 & 0.06 \\
 & DINOv2-S LSH-256 & 0.62 & 0.51 & 0.47 & 0.66 & 0.21 \\
 & ISC21 1st\reftier & 0.62 & 0.76 & 0.60 & 0.78 & 0.09 \\
 & ISCC Image-64 & 0.62 & 0.47 & 0.43 & 0.66 & 0.62 \\
 & PDQ & 0.62 & 0.49 & 0.47 & 0.64 & 0.00 \\
 & SSCD R50\reftier & 0.62 & 0.68 & 0.58 & 0.72 & 0.18 \\
\addlinespace
AudioSeal & CLAP & 0.24 & 0.37 & 0.12 & 0.49 & -- \\
 & Chromaprint & 0.24 & 0.67 & 0.21 & 0.71 & -- \\
 & ISCC Audio-64 & 0.24 & 0.40 & 0.16 & 0.48 & -- \\
 & audfprint & 0.24 & 0.64 & 0.23 & 0.65 & -- \\
\addlinespace
Video Seal & ISCC Video-64 & 0.92 & 0.62 & 0.62 & 0.92 & -- \\
 & SSCD frames\reftier & 0.92 & 0.54 & 0.54 & 0.92 & -- \\
 & TMK+PDQF & 0.92 & 0.64 & 0.64 & 0.92 & -- \\
 & vPDQ & 0.92 & 0.58 & 0.58 & 0.92 & -- \\
\bottomrule
\end{tabular}

\end{table}


\begin{figure*}[!t]
\centering
\includegraphics[width=\textwidth]{fp_fig11_wm_complementarity}
\caption{Share of attacked copies satisfying only the watermark criterion (blue), only the fingerprint criterion (orange), or both (grey), per transformation. Watermark success means exact-key recovery for images and audio, and expected-key verification for video; the video panel does not measure unique registry identification. (a) Images: TrustMark-Q and DINOv2-S, with the fingerprint threshold calibrated to pair-level FPR $10^{-7}$. (b) Audio: AudioSeal and Chromaprint. (c) Video: Video Seal and the ISCC Video-Code. Audio and video fingerprint thresholds are calibrated to query-level FPR 1\%; achieved held-out rates are reported separately.}
% The figure labels use the method-specific success definitions in the caption.
\label{fig:fp-wm}
\end{figure*}
